\scanpage{72}

\deleted{\diagram{\begin{tikzcd}[column sep=small]
\deleted{C^*(X,\Omega^*)}\arrow[r]&
\deleted{\mathbb R^*\Gamma(X,\Omega^*)}&
\mathbb R^*\Gamma(X,k)\arrow[l,"\varepsilon_{\Omega^*}"']\arrow[r,"\sim","{\varepsilon_{S_X^*}}"']\arrow[d,"\sim"']\arrow[dr,dashed,"\varepsilon_{S_X^{\prime *}}"]&
\mathbb R^*\Gamma(X,S_X^*)\arrow[d,"\sim"]&
\Gamma(X,S^*_{X,\unsure{k}})=C^*(S(X)_*,k)\arrow[l,"\sim"']\arrow[d,"\sim"]\\
&&\mathbb R^*\Gamma(X,\Omega^*)\arrow[r,"\sim"]&
\mathbb R^*\Gamma(X,S_X^{\prime *})&
\Gamma(X,S_X^{\prime *})=C^*(S(X)'_*,k)\arrow[l,"\sim"']\\
&&C^*(X,\Omega^*)\arrow[u]&&
\end{tikzcd}}}
\margin{($S_X^*$ \unsure{homot$^t$ fin}; \unsure{système paracompact})}

\uline{NB} \unsure{les hyp.} (4) \unsure{jointes|jointe} avec
$\unsure{K_*|k_*}$ définit un hom
$\Omega^*\longrightarrow S_X^{\prime *}$ compatible avec les augmentations,
d'où carré comm.
\diagram{\begin{tikzcd}
\Omega^*\arrow[r,"\unsure{\text{par }K}"]&S_X^{\prime *}\\
k\arrow[u]\arrow[r,dash]&S_X^*\arrow[u]
\end{tikzcd}}
\deleted{\unsure{(isom.}}
(d'isom dans $D(X,S)$) donnant lieu à un cube commutatif

\diagram{\begin{tikzcd}[column sep=small,row sep=large]
C^*(X,\Omega^*)\arrow[rr]\arrow[dd,"\unsure{q}"]\arrow[dr,phantom,"\deleted{\unsure{\searrow}}"]&&
C^*(X,S_X^{\prime *})\overset{\mathrm{dfn}}=C(S'_X,k)\arrow[dd,"\sim"]&\\
&[C(X,k)]\arrow[ul,dashed]\arrow[rr,dashed]\arrow[dd,dashed]&&
C^*(X,S_X^*)=C(S_X,k)\arrow[ul,"\sim\ \text{pr. sur }S'_X"']\arrow[dd,"\sim\ (\text{Cartan, }S_X^*\text{ est }\unsure{\text{``homot}^t\text{ fin''}})"]\\
\mathbb R^*\Gamma(X,\Omega^*)\arrow[rr,"\sim"]&&\mathbb R^*\Gamma(X,S_X^{\prime *})&\\
&\mathbb R^*\Gamma(X,k)\arrow[ul,"\sim"]&&
\mathbb R^*\Gamma(X,S_X^*)\arrow[ul,"\sim"]\arrow[ll,"\unsure{\sim;\ \leftarrow|\rightarrow}"]
\end{tikzcd}}

\margin{\deleted{$\longleftarrow\mathbb R^*\Gamma(X,S^*_{X,\unsure{k}})$}}

Le bas du cube est formé d'isom, car transformé par $\mathbb R^*\Gamma$ d'un carré d'iso dans une cat. dérivée; la face \unsure{droite} est formée d'iso, car on le sait déjà pour deux arêtes, donc aussi pour la quatrième qui reste.
\margin{71}
